\chapter*{مقدمه}
\addcontentsline{toc}{chapter}{مقدمه}
\markboth{مقدمه}{مقدمه}
\label{chap:introduction}
خودکشی یکی از چالش‌های جدی سلامت عمومی است و پیامد آن تنها به یک فرد محدود نمی‌شود؛ خانواده، شبکۀ اجتماعی و نظام سلامت نیز از آن تأثیر می‌پذیرند. سازمان جهانی بهداشت بر ضرورت پیشگیری، شناسایی نشانه‌های هشدار و دسترسی به حمایت مناسب تأکید می‌کند \cite{ref1}. تعامل عوامل روان‌شناختی، اجتماعی، زیستی و محیطی، پیش‌بینی این پدیده را دشوار می‌سازد \cite{ref3}. ازاین‌رو هر ابزار محاسباتی باید در کنار این پیچیدگی و محدودیت‌های داده تفسیر شود و جایگزین ارزیابی متخصص تلقی نشود.

گسترش شبکه‌های اجتماعی حجم بزرگی از متن‌های خوداظهاری را در دسترس پژوهشگران قرار داده است. انتخاب واژه، اشاره به تنهایی و ناامیدی، تغییر لحن و شیوۀ روایت تجربۀ فردی می‌توانند با وضعیت روانی ارتباط داشته باشند \cite{ref6,ref7,ref8,ref9}. پردازش زبان طبیعی و یادگیری ماشین راهی برای تحلیل نظام‌مند این الگوها فراهم می‌کنند. با این حال، متن شبکۀ اجتماعی معادل پروندۀ بالینی یا تشخیص قطعی نیست. منبع برچسب، زمان گردآوری، زبان، فرهنگ و تفاوت کاربران می‌توانند اعتبار و تعمیم نتیجه را محدود کنند \cite{ref10,ref11}.

در حوزۀ سلامت روان، حریم خصوصی یک پیش‌شرط اساسی است. متن ممکن است هم‌زمان شامل اطلاعات هویتی، تجربۀ شخصی، وضعیت درمان و روابط اجتماعی باشد. تمرکز چنین داده‌ای در یک مخزن، سطح مواجهه با خطای دسترسی، افشا و استفادۀ ثانویه را افزایش می‌دهد. دغدغه‌های اخلاقی و قواعدی مانند \en{GDPR} و \en{HIPAA} نیز اهمیت تعیین مسئولیت، رضایت، محدودیت استفاده و حفاظت از داده را نشان می‌دهند \cite{ref4}. بنابراین افزایش یک معیار پیش‌بینی، بدون توجه به چرخۀ داده، پاسخ کاملی به مسئله نیست.

یادگیری مشارکتی، با محلی نگه داشتن دادۀ آموزش و اشتراک مؤلفه‌های مدل یا خروجی‌ها، امکان همکاری میان منابع داده را بررسی می‌کند \cite{ref15,ref16,ref17}. این ویژگی می‌تواند نیاز به انتقال متن خام را کاهش دهد، اما به‌خودی‌خود نشت اطلاعات را منتفی نمی‌کند. به‌روزرسانی مدل یا احتمال پیش‌بینی نیز ممکن است حامل اطلاعات حساس باشد. همچنین هویت فرستنده، تمامیت پیام، بازپخش، اعتبار وزن و رفتار عضو مجاز اما مخرب، مسائلی جداگانه‌اند که صرف توزیع آموزش آن‌ها را حل نمی‌کند.

رمزگذاری همریخت اجازه می‌دهد برخی عملیات، از جمله جمع وزن‌دار، در دامنۀ متن رمز انجام شوند \cite{ref19,ref20,ref21,ref22}. اگر احتمال‌های خروجی کلاینت‌ها به نمایش صحیح معتبر تبدیل شوند، یک طرح جمعی می‌تواند متن رمزِ مجموع را بدون رمزگشایی جداگانۀ همۀ مشارکت‌ها بسازد. امضای دیجیتال نیز می‌تواند اصالت فیلدهای محموله را پیش از پذیرش کنترل کند. با این حال، انجام این دو کار باید با مدل اعتماد روشن همراه باشد: سروری که کلید خصوصی را دارد، صرفاً با انتخاب مسیر برنامۀ «فقط تجمیع» از توان رمزگشایی منفرد محروم نمی‌شود.

\section*{انگیزه و هدف پژوهش}
انگیزۀ اصلی، بررسی هم‌زمان کیفیت تحلیل متن و الزامات حفاظت از تبادل خروجی‌هاست. هدف، توسعۀ یک ارزیابی قابل‌بازتولید برای انتخاب طبقه‌بند، بررسی اثر تعداد کلاینت‌ها بر خروجی تجمیعی و صورت‌بندی یک لایۀ احراز اصالت سازگار با جبر جمعی است. این اهداف به سه آزمایش و تحلیل مرتبط، اما متمایز منجر می‌شوند:
\begin{enumerate}
\item مقایسۀ متمرکزِ خط‌لوله‌های متن با جداسازی آموزش، اعتبارسنجی و آزمون و جلوگیری از یادگیری واژگان از بخش نگه‌داشته‌شده؛
\item آموزش مدل‌های مستقل محلی با تخصیص طبقه‌بندی‌شدۀ همۀ ردیف‌های آموزش در سناریوهای ۱۰، ۲۰، ۵۰ و ۱۰۰ کلاینت و تجمیع احتمال بر اساس اندازۀ دادۀ محلی؛
\item تحلیل و آزمون کارکردیِ تجمیع مجموع صحیح وزن‌دار با بنالوح و بررسی امضای \en{Ed25519} پیش از پذیرش محموله.
\end{enumerate}

\section*{پرسش‌ها و فرض‌های قابل بررسی}
پرسش نخست این است که با بازنمایی \en{Word2Vec} و ارزیابی بدون نشت، چهار خانوادۀ جنگل تصادفی، درخت تصمیم، \en{XGBoost} و رأی‌گیری بیز ساده چه عملکردی بر اعتبارسنجی دارند و انتخاب تثبیت‌شده چه نتیجه‌ای در آزمون می‌دهد. برتری هیچ مدل از پیش به‌عنوان نتیجه فرض نمی‌شود. معیار انتخاب، $F_1$ کلاس مثبت و قواعد رفع تساوی تعریف‌شده است؛ دقت آزمون برای انتخاب مدل استفاده نمی‌شود.

پرسش دوم به اثر اندازۀ دادۀ محلی مربوط است. هنگامی که کل داده ثابت و نسبت کلاس‌ها تقریباً حفظ شود، افزایش تعداد کلاینت چه تغییری در واژگان محلی، پوشش متن و کیفیت احتمال تجمیعی ایجاد می‌کند؟ این پرسش با یک تقسیم طبقه‌بندی‌شدۀ تقریباً هم‌توزیع بررسی می‌شود. نتیجۀ آن نباید به سناریوهای غیرهم‌توزیع، تقسیم زمانی یا استقلال واقعی کاربران تعمیم داده شود.

پرسش سوم این است که چگونه می‌توان هویت، تمامیت، دور و یکتایی محموله را پیش از ضرب متن‌های رمز کنترل کرد، بدون آنکه امضا در ساختار جبری بنالوح دخالت کند. درستی مجموع به شرط ظرفیت فضای پیام وابسته است؛ امنیت محموله نیز به ثبت معتبر کلید، محرمانه ماندن کلید امضا و نگهداری وضعیت بازپخش وابسته است. این پیش‌شرط‌ها جزئی از نتیجه‌اند، نه جزئیاتی قابل حذف از گزارش.

\section*{روش و اصلاحات ارزیابی}
منبع داده، مجموعۀ انگلیسی \en{Suicide Watch} با ۲۳۲٬۰۷۴ ردیف خام است. پس از حذف تکرار خام، ۲۳۱٬۹۸۱ ردیف تقسیم می‌شود. سهم‌های اسمی آموزش، اعتبارسنجی و آزمون ۶۴، ۱۶ و ۲۰ درصدند. پاکسازی و حذف هم‌پوشانیِ ناشی از آن، تعدادهای نهایی را به ۱۴۸٬۳۸۰، ۳۷٬۰۸۶ و ۴۶٬۳۴۲ می‌رساند. این شمارش جایگزین توصیف قبلیِ ۴۶۴٬۱۴۸ نمونه است.

پاکسازی، نفی انگلیسی، اعداد و حروف یونیکد را حفظ می‌کند و فهرست آموخته‌شده‌ای از کلمات توقف ندارد. بازنمایی صدبعدیِ میانگین بردار واژه درون خط‌لوله قرار دارد و در هر تقسیم اعتبارسنجی متقابل فقط روی آموزش همان تقسیم یاد گرفته می‌شود. مدل و آستانه بر اعتبارسنجی جداگانه تثبیت می‌شوند و آزمون نهایی برای آن انتخاب گزارش می‌شود. ابزار فارسی «دنا» به‌عنوان دستاورد جانبی معرفی می‌شود، اما نتایج این مسیر انگلیسی به آن نسبت داده نمی‌شوند.

در مسیر مشارکتی، هر کلاینت هم بازنمایی و هم طبقه‌بند خود را محلی یاد می‌گیرد. میانگین وزن‌دار احتمال‌ها محاسبه می‌شود؛ وزن درخت‌ها یا مختصات بازنماییِ مستقل میانگین‌گیری نمی‌شوند. به این ترتیب، روش اجراشده با آموزش چنددوره‌ای \en{FedAvg} تفاوت دارد. در اجرای ثبت‌شده، سناریوی ۱۰ کلاینت بر اساس $F_1$ اعتبارسنجی انتخاب می‌شود و دقت آزمون آن با آستانۀ منتخب ۹۱٫۵۸ درصد است. فقط یک بذر برای هر سناریو وجود دارد؛ عدم‌قطعیت بین اجراها از این داده قابل برآورد نیست.

در مسیر رمزنگاری، طرح‌های جمعی و ضربی ابتدا از نظر جبری تفکیک می‌شوند. سپس \en{SBHE-Benaloh} مقدارهای چهاررقم‌اعشاری را به صحیح تبدیل، محمولۀ دارای وزن و شناسۀ دور را امضا و پیش از تجمیع اعتبارسنجی می‌کند. نتایج طبقه‌بندی مشارکتی از تجمیع آشکارِ اعشاری و زمان‌های رمزنگاری از بنچمارک‌های عددی جداگانه به دست آمده‌اند؛ سامانۀ رمزگذاری‌شدۀ سرتاسریِ طبقه‌بندی در این دو آزمایش اثبات یا اندازه‌گیری نشده است.

\section*{سهم پژوهش و محدودیت‌ها}
سهم اصلی، کنار هم قرار دادن ارزیابیِ قابل ممیزی، تعریف صریح توزیع دادۀ کلاینت‌ها و تفکیک نقش احراز اصالت از همریختی است. نتیجۀ مدل‌ها به‌جای اتکا به ارزیابیِ دارای نشت یا مقایسۀ مجموعه‌های متفاوت، با پروتکل مشخص گزارش می‌شود. نتیجۀ رمزنگاری نیز به‌جای ادعای محرمانگی مطلق، با شرط ظرفیت و مرز اعتماد سرور بیان می‌شود. معرفی ابزار دنا و مستندسازی خروجی‌ها، بازاستفادۀ پژوهشی را تسهیل می‌کند.

دادۀ غیربالینی و انگلیسی، یک بذر در آزمایش مشارکتی، تعداد محدود پیکربندی‌های رمزنگاری و پارامترهای کوچکِ پژوهشی، محدودیت‌های اصلی‌اند. در کلاینت‌های واقعی، ناهمگونی زبان و موضوع، تغییر توزیع زمانی، عدم توازن شدید و نمونه‌های چندگانۀ یک کاربر می‌توانند نتیجه را تغییر دهند. امضای معتبر نیز راست‌گویی یک عضو مجاز را ثابت نمی‌کند. اثبات بازه، رمزگشایی آستانه‌ای، حریم خصوصی تفاضلی و اثبات دانش صفر در پیاده‌سازی حاضر اجرا نشده‌اند.

\section*{ساختار پایان‌نامه}
فصل اول ضرورت، روش کلی و پیشینۀ سه حوزۀ تحلیل متن، یادگیری مشارکتی و محاسبۀ همریخت را مرور می‌کند. فصل دوم مبانی پردازش زبان طبیعی، طراحی اصلاح‌شدۀ ارزیابی، بازنمایی متن، طبقه‌بندها و نتایج متمرکز و مشارکتی را ارائه می‌دهد. فصل سوم به حفاظت از ورودی، تسهیم راز، فرمول‌های پنج خانوادۀ رمزگذاری و محدودیت مقایسۀ زمانی آن‌ها اختصاص دارد. فصل چهارم ساختار محموله، درستی و حدود امنیت \en{SBHE-Benaloh} و نتایج زمانی ذخیره‌شده را تشریح می‌کند. در پایان، پیوست‌های دنا، لگاریتم گسسته، دیفی\lr{--}هلمن، باقیماندۀ مرکب و \en{Ed25519}، همراه با اختصارات، واژه‌نامه و منابع آمده‌اند.
